%%=============================================================================
%% Realisatie
%%=============================================================================
\chapter{\IfLanguageName{dutch}{Realisatie}{Implementation}}%
\label{ch:realisatie}
Dit hoofdstuk beschrijft hoe de functionele en technische vereisten van het
graduaatsproject worden vertaald naar een zelfstandige proof of concept (PoC)
van een visuele e-maileditor. De nadruk ligt op de belangrijkste onderdelen
van de oplossing, de gemaakte technische keuzes en de aandachtspunten bij de
verwerking van e-mailtemplates.

De realisatie wordt besproken aan de hand van de codeopbouw, de werkwijze en
het versiebeheer, de belangrijkste technische onderdelen en de problemen die
tijdens de ontwikkeling worden onderzocht. De concrete implementatie en de
uiteindelijke resultaten worden beschreven op basis van de effectief
gerealiseerde functionaliteiten.

\section{\IfLanguageName{dutch}{Opbouw van de code}{Structure of the code}}%
\label{sec:opbouw-code}
De proof of concept wordt opgezet als een zelfstandige toepassing, los van de
productieomgeving van Vesalius.ai. Hierdoor kunnen de editor, de verwerking
van e-mailtemplates en de gegenereerde uitvoer afzonderlijk worden ontwikkeld
en getest, zonder de bestaande e-mailverwerking in productie te wijzigen.

De beoogde architectuur bestaat uit een frontend, een mock-API en een
verwerkingslaag voor e-mailtemplates. De frontend wordt ontwikkeld met
Angular en TypeScript. De mock-API wordt opgezet met Node.js, TypeScript en
Express en staat in voor de communicatie tussen de frontend en de logica voor
templatebeheer.

Binnen de frontend bevindt zich de gebruikersinterface waarmee een gebruiker
een e-mailtemplate kan selecteren, de inhoud visueel kan aanpassen en het
resultaat vooraf kan bekijken. Voor de WYSIWYG-functionaliteit wordt Quill 2
als kandidaat onderzocht. De definitieve keuze wordt bepaald door de
technische haalbaarheid en de ondersteuning van de vereiste
bewerkingsmogelijkheden.

De backendlaag behandelt het laden en opslaan van templates en coördineert
de verwerking van de inhoud. De opslag wordt zo opgezet dat de PoC niet
onnodig afhankelijk wordt van één specifieke databasetechnologie. Afhankelijk
van de implementatie kunnen hiervoor tijdelijke opslag of vooraf gedefinieerde
testgegevens worden gebruikt.

De renderingpipeline vormt een afzonderlijk onderdeel van de toepassing.
Deze laag verwerkt de HTML-inhoud, controleert toegestane placeholders en
combineert de bewerkbare inhoud met de vaste structuur van de e-mail. Indien
nodig wordt de CSS inline geplaatst om de compatibiliteit met e-mailclients
te verbeteren.

De lokale testomgeving ondersteunt het controleren van de gegenereerde
e-mails. Hiervoor wordt Mailpit als mogelijke testmailservice gebruikt.
Docker kan worden ingezet om de benodigde services reproduceerbaar op te
starten.

Door de frontend, de API, de templateopslag en de renderingpipeline logisch
van elkaar te scheiden, blijven de verantwoordelijkheden duidelijk. Dit
vereenvoudigt het afzonderlijk testen van de onderdelen en maakt het
eenvoudiger om de PoC later technisch te evalueren met het oog op een
mogelijke integratie in Vesalius.ai.

\section{\IfLanguageName{dutch}{Werkwijze en versiebeheer}{Way of working and version control}}%
\label{sec:werkwijze}
De werkzaamheden worden opgevolgd via Jira. De beschikbare tickets en
projectvereisten vormen het uitgangspunt voor de planning en de uitvoering
van de werkzaamheden. Per onderdeel wordt bepaald welke functionaliteit nodig
is, welke technische aandachtspunten onderzocht moeten worden en hoe het
resultaat kan worden gecontroleerd.

De broncode wordt beheerd met Git en Bitbucket binnen de interne
ontwikkelomgeving van Vesalius.ai. Hierdoor kunnen wijzigingen worden
opgevolgd en kan de evolutie van de implementatie worden nagegaan. De
repository is intern en wordt niet als publieke broncode beschikbaar gesteld.

Voor de ontwikkeling wordt Visual Studio Code gebruikt. De lokale
ontwikkelomgeving draait op Windows 11 met Ubuntu via WSL2. Docker wordt
gebruikt binnen de ontwikkelworkflow wanneer containers nodig zijn voor de
testomgeving of ondersteunende services.

De ontwikkeling volgt een iteratieve werkwijze. Eerst worden de vereisten en
de technische haalbaarheid onderzocht. Vervolgens wordt de basisstructuur van
de toepassing opgezet, waarna de editor, het templatebeheer en de
renderingpipeline stapsgewijs worden uitgewerkt. Daarna worden de
previewfunctionaliteit, de testmailomgeving en de geautomatiseerde tests
toegevoegd volgens de projectplanning.

Na iedere relevante wijziging wordt gecontroleerd of de aangepaste
functionaliteit aan de verwachte vereisten voldoet. Daarbij wordt niet alleen
rekening gehouden met de zichtbare werking van de editor, maar ook met de
HTML-uitvoer, de verwerking van placeholders en de veiligheid van de
gegenereerde e-mail.

De precieze voortgang, de uitgevoerde taken en eventuele afwijkingen worden
bijgehouden aan de hand van de projectplanning en de bijbehorende tickets.
De definitieve werkwijze en de werkelijk uitgevoerde controles worden in dit
hoofdstuk aangevuld op basis van de gerealiseerde implementatie.

\section{\IfLanguageName{dutch}{Belangrijkste onderdelen}{Key components}}%
\label{sec:onderdelen}
De realisatie is opgebouwd rond drie belangrijke technische onderdelen:
de visuele editor en het templatebeheer, de renderingpipeline en de
testomgeving voor de gegenereerde e-mails.

\subsection{Visuele editor en templatebeheer}
De visuele editor vormt het centrale onderdeel van de PoC. Het doel is om
gebruikers toe te laten de inhoud en ondersteunde opmaak van vooraf bepaalde
e-mailtemplates aan te passen zonder rechtstreeks HTML-code te moeten
bewerken.

De gebruiker selecteert een van de drie piloottypes: een
afspraakbevestiging, een afspraakherinnering of een afspraakannulering.
Vervolgens wordt de bijbehorende template geladen in de editor. De gebruiker
kan de beschikbare inhoud aanpassen en het resultaat vooraf bekijken.

De editor richt zich op de inhoud van de e-mail. De vaste e-mailschil, met
bijvoorbeeld het logo, de merkkleuren en de footer, blijft grotendeels
buiten de bewerkingsmogelijkheden. Hierdoor blijft de scope van de PoC
beperkt en kan de aandacht uitgaan naar het aanpassen en betrouwbaar
renderen van de e-mailinhoud.

Naast het bewerken is ook het opslaan en opnieuw laden van templates van
belang. De opslaglaag wordt daarom gescheiden van de editorcomponent.
Hierdoor hoeft de interface niet rechtstreeks afhankelijk te zijn van de
gekozen opslagtechnologie.

De concrete implementatie wordt geëvalueerd op basis van het laden van
templates, het aanpassen van inhoud, het bewaren van wijzigingen en het
opnieuw openen van een opgeslagen template.

\subsection{Renderingpipeline en veilige verwerking}
Een tweede belangrijk onderdeel is de omzetting van de bewerkte template
naar uiteindelijke e-mailinhoud. De HTML die door een visuele editor wordt
gegenereerd, kan niet zonder controle worden gebruikt. De invoer moet worden
gevalideerd en niet-toegestane HTML-elementen of attributen moeten worden
verwijderd of geweigerd.

De renderingpipeline verwerkt daarom de inhoud volgens vooraf bepaalde
regels. Eerst wordt de HTML gesanitiseerd. Vervolgens worden uitsluitend
toegestane placeholders verwerkt met fictieve testgegevens. Daarna wordt de
inhoud gecombineerd met de vaste e-mailschil. Indien nodig wordt CSS inline
geplaatst om de weergave in e-mailclients te ondersteunen.

Bij de placeholderverwerking is het belangrijk dat onbekende of ontbrekende
waarden veilig worden afgehandeld. Ook moet rekening worden gehouden met de
context waarin een waarde wordt ingevoegd. Het ontsnappen van waarden en het
controleren van toegestane invoer blijven noodzakelijk om ongewenste HTML-
of scriptinhoud te vermijden.

De renderingpipeline wordt afzonderlijk beoordeeld van de editor. Een
correct werkende editor garandeert immers niet dat de uiteindelijke HTML
veilig is of in iedere e-mailclient hetzelfde wordt weergegeven.

\subsection{Preview, testmails en compatibiliteitscontroles}
Het derde onderdeel omvat de controle van de gegenereerde e-mail. De
previewfunctie maakt het mogelijk om de gerenderde inhoud te beoordelen
voordat een testmail wordt verstuurd. Daarbij wordt nagegaan of de inhoud,
de placeholders en de vaste e-mailschil correct worden weergegeven.

Voor het testen van e-mails wordt een lokale testomgeving voorzien. Mailpit
kan worden gebruikt om testmails op te vangen en te inspecteren zonder dat
deze naar echte ontvangers worden verstuurd. De testgegevens zijn fictief
en bevatten geen echte patiëntinformatie.

Daarnaast worden geautomatiseerde tests voorzien voor onderdelen zoals
templatevalidatie, placeholderverwerking en HTML-sanitatie. Vitest kan
worden ingezet voor unit- en integratietests en Playwright voor
browsergebaseerde tests, afhankelijk van de uiteindelijke implementatie.

Ten slotte wordt de weergave van de gegenereerde e-mails onderzocht in
beschikbare e-mailclients, waaronder Gmail, Outlook en Apple Mail. Eventuele
verschillen in HTML- en CSS-ondersteuning worden geregistreerd. De resultaten
van deze controles worden opgenomen in de bijlage met testresultaten.

\subsection{Broncode en vertrouwelijkheid}
De broncode van de toepassing bevindt zich in de interne Bitbucket-omgeving
van Vesalius.ai. Omdat de repository deel uitmaakt van de bedrijfsomgeving,
worden interne bronbestanden, bedrijfsgegevens en vertrouwelijke configuratie
niet zonder toestemming opgenomen in dit rapport.

De technische werking kan worden toegelicht met architectuurschema's,
beschrijvingen van de gegevensstroom en vereenvoudigde codefragmenten.
Wanneer een codefragment uit de interne implementatie niet mag worden
gepubliceerd, wordt de relevante logica op een hoger niveau beschreven.
Zo blijft de technische onderbouwing van het project mogelijk zonder
vertrouwelijke bedrijfsinformatie prijs te geven.

\section{\IfLanguageName{dutch}{Problemen onderweg}{Problems along the way}}%
\label{sec:problemen}
Een belangrijk aandachtspunt bij de realisatie is de betrouwbaarheid van de
gegenereerde e-mail. Een WYSIWYG-editor kan inhoud visueel correct weergeven,
maar de resulterende HTML moet ook veilig worden verwerkt en bruikbaar zijn
in verschillende e-mailclients. Dit vraagt om een afzonderlijke controle van
de editoruitvoer en de uiteindelijke rendering.

Een tweede aandachtspunt is het correct verwerken van placeholders. Dynamische
waarden moeten gecontroleerd worden ingevuld en mogen niet leiden tot
onverwachte HTML-uitvoer. Daarom moet vooraf worden bepaald welke placeholders
toegestaan zijn en hoe ontbrekende of onbekende waarden worden afgehandeld.

Daarnaast is de ondersteuning van HTML en CSS niet identiek in alle
e-mailclients. Een e-mail die in de browserpreview correct wordt weergegeven,
kan in Outlook, Gmail of Apple Mail afwijkingen vertonen. De preview alleen
volstaat daarom niet als bewijs van compatibiliteit. Waar mogelijk worden
gegenereerde testmails ook in de beschikbare e-mailclients gecontroleerd.

Ook de keuze van de opslagmethode vraagt aandacht. Omdat het project een
zelfstandige PoC is en geen productie-integratie, moet de opslag voldoen aan
de testvereisten zonder onnodige afhankelijkheden te introduceren. De
opslaglaag wordt daarom zo veel mogelijk losgekoppeld van de rest van de
toepassing.

Tot slot moet de scope van het project bewaakt worden. De PoC is bedoeld om
de haalbaarheid van een visuele e-maileditor aan te tonen en niet om de
bestaande e-mailinfrastructuur van Vesalius.ai volledig te vervangen.
Productieverzending, integratie met de bestaande platformarchitectuur en
verwerking van echte patiëntgegevens vallen buiten de scope.

De uiteindelijke bespreking van problemen wordt aangevuld met de concrete
technische obstakels die tijdens de implementatie werden vastgesteld, de
oplossingen die daarvoor werden gekozen en de resultaten van de uitgevoerde
hertests.
